\section{Reconstruction of the retained physical anchor input}
\label{sec:sc-geometry}
This appendix reconstructs the specific consequences of native modules 287 and 294 needed in the new sublevel proof. Its statements are inherited results, not newly claimed contributions. The exact remote body of 287 is retained under \texttt{provenance/287-original.tex}. The full newer 294 manuscript remains in the pinned complete local predecessor archive.

\begin{sclemma}[Buffered all-direction line covering, inherited]\label{lem:sc-buffer}
Every straight segment of length $H=28/5$ has a lattice center at perpendicular distance at most $\omega=179/400$ from its line, with perpendicular foot strictly inside the segment.
\end{sclemma}
\begin{proof}
For a primitive triangular lattice vector of length $l$, the parallel lattice rows have spacing $h=\sqrt3/(2l)$. At a segment midpoint $q$ choose a nearest row and then a nearest lattice point on it. The two coordinates of $z-q$ parallel and perpendicular to this primitive direction obey $|x|\le l/2$, $|y|\le h/2$. For a direction at angle $\delta$ to that vector, its transverse projection is at most
$W=(h|\cos\delta|+l|\sin\delta|)/2$, and its longitudinal projection at most $\sqrt{l^2+h^2}/2$.

By lattice symmetries reduce to an unoriented angle $0\le\theta\le\pi/6$ and put $t=\tan\theta/\sqrt3$. The primitive lattice vectors and the bounds obtained from their row rectangles are
\[
\begin{array}{c|c|c|c}
 t\text{ interval}&v&l^2&\sup W^2\\\hline
 [0,1/60]&(1,0)&1&961/4804\\
 [1/60,1/12]&(11/2,\sqrt3/2)&31&13697401/73866304\\
 [1/12,1/5]&(5/2,\sqrt3/2)&7&529/3136\\
 [1/5,1/3]&(3/2,\sqrt3/2)&3&75/448
\end{array}
\]
Here is an explicit check that these are bounds for whole intervals. For the first row,
$W^2=3(1+2t)^2/(16(1+3t^2))$, increasing since $2-3t>0$. For each other row write $v=((2a+1)/2,\sqrt3/2)$ with $a=5,2,1$, $s=2a+1$, $N=a^2+a+1$. Then
\[
 W=\frac{\sqrt3}{4\sqrt{1+3t^2}}
       \left(\frac{s+3t}{2N}+|st-1|\right).
\]
On the listed interval $t\le1/s$, the positive numerator decreases with derivative $3/(2N)-s<0$, and the denominator increases. Evaluate its left endpoint, and the right endpoint for the first row, to obtain the table. Every listed rational is strictly below $(179/400)^2$. The largest longitudinal bound is
$\sqrt{31+3/124}/2=\sqrt{3847/124}/2<H/2$, since $96175<97216$. Therefore the selected foot lies strictly inside the original segment. This argument neither changes the disk radii nor discards a direction.
\end{proof}

\begin{sclemma}[Physical-time turning and the 94-flight anchor, inherited]\label{lem:sc-anchor}
Every actual regular trajectory segment of physical duration $H$ has total interior incidence cosine at least $\gamma=1/4480$. Every actual free flight has length at most $H$, and every complete 94-flight block has an interior collision with $c_i\ge1/416640$.
\end{sclemma}
\begin{proof}
A collision changes the unit velocity by a vector of length $2c_i$. Hence for a physical segment $q(s)$, $0\le s\le H$, with initial outgoing velocity $v$, direct integration gives
\[
 |q(s)-q(0)-sv|\le2H\sum_{0<t_i<H}c_i.
\]
Endpoint collisions do not affect this estimate: an initial collision is accounted for by the outgoing $v$, and the velocity change at $H$ does not affect positions up to $H$. If the sum were less than $1/4480$, the deviation would be less than $1/400$. Apply Lemma~\ref{lem:sc-buffer} to the initial straight reference segment and evaluate at its interior perpendicular foot. The actual position would then have distance less than
$179/400+1/400=9/20\le R$ from a lattice center, contradicting exterior nonpenetration. This proves the incidence sum.

A straight free flight of length exceeding $H$ would contain an $H$ segment intersecting a disk interior by the same covering lemma, a contradiction. Consecutive distinct disks are separated by $\tau_-=3/50$, so 94 flights take at least $94\tau_-=141/25>H$. Apply the physical-time bound to their first $H$ units. The relevant collisions are among the 93 interior collisions of the block. Thus their total cosine is at least $1/4480$, and one is at least $1/(4480\cdot93)=1/416640$. All collisions are actual; no probabilistic independence, symbolic model or number of words is used.
\end{proof}
